\documentclass[11pt,a4paper]{article}
\usepackage[utf8]{inputenc}
\usepackage[T1]{fontenc}
\usepackage{amsmath,amssymb}
\usepackage{booktabs}
\usepackage{geometry}
\usepackage{hyperref}
\usepackage{xcolor}
\usepackage{graphicx}
\usepackage{xurl}
\usepackage{enumitem}
\geometry{margin=25mm}
\hypersetup{colorlinks=true,linkcolor=blue!60!black,urlcolor=blue!60!black}

\title{Lift Model Packs: Exploratory Modelling Reference\\[4pt]
\large Research Reference for Epic \#11740 (LIFT-11, Issue \#11751)}
\author{UpstreamDrift Research Reference}
\date{Revision 2026-10-09}

\begin{document}
\maketitle

\begin{abstract}
Four per-engine model packs (\texttt{OpenSim\_Models}, \texttt{MuJoCo\_Models},
\texttt{Drake\_Models}, \texttt{Pinocchio\_Models}) each build full-body models
of the five classical barbell lifts --- back squat, deadlift, bench press,
snatch, and clean and jerk --- for a reference lifter with a competition
barbell. A cross-engine audit (epic \#11740, LIFT-1 \#11741) reduces each
pack's own forward kinematics to bar, hand and foot geometry and compares the
four packs pose by pose. This reference collects the exploratory modelling
material around that audit that is not itself the governed calculation:
the model inventory per pack, the phase and trajectory conventions and why a
shared standard remains open, what inverse-dynamics and spine-load analysis
exists today (Movement Optimizer, not the engine packs), the measured parity
numbers, the derived discrepancy list, and the documented negative results.
\end{abstract}

\noindent\textbf{Document status.} This is an editable standalone research
source, separate from the engineering design manual. The sole editable
design-manual source remains \texttt{manuals/upstreamdrift} (QMD); the
governed calculation for this material is Chapter 15, ``Lift Model Pack
Audit'' (\texttt{manuals/upstreamdrift/chapters/15-lift-pack-audit.qmd}),
itself recorded there as evidence class ``derived, provisional'' and not yet
approved. This document is a software and model-consistency check, not
scientific validation against capture data: no lift motion-capture or
force-plate data exist for any of the five lifts, and no uncertainty against
a measured lift is claimed anywhere below.

\section{Purpose and Scope}
The four packs are parametric Python generators that build a
model from body mass, height and plate mass. The audit compares them on the
same canonical inputs: it checks masses and joint limits against a shared
biomechanics parity standard, reduces each engine's forward kinematics to
bar/hand/foot geometry, and derives a discrepancy list with tracking issues
(Chapter 15, ``Purpose''). What the audit explicitly does \emph{not} do is
restated here because it bounds every claim in this document: \emph{``It
measures software parity between the packs only. It does not simulate a
lift, and it does not compute joint moments, spine loads or muscle
forces.''} (Chapter 15, ``Purpose''). This reference does not re-derive the
audit's geometry equations (bar frame, hand-to-bar offsets, placement
independent pose summary, same-input parity, lifter-only centre of mass,
segment mass deviation, joint-limit comparison, grip width, start pose,
reference FK) --- those ten equations and their tolerances are Chapter 15's
governed content and are only cited here by name. The user-facing summary of
the same material lives in \texttt{docs/help/movement\_optimizer.md}
(``Lifting Models: Modelling Reference'' section), which this document
extends with the measured numbers and the open-gap bookkeeping.

\section{Model Inventory per Pack}

\subsection{Body Model}
All four packs scale a full-body model from \texttt{body\_mass} and
\texttt{height} to Winter (2009) anthropometrics, declared as
\texttt{anthropometrics: winter\_2009} in each pack's \texttt{model\_pack.yaml}.
OpenSim's is a 15-segment body (pelvis, torso, head, bilateral arms and legs)
with ball joints at hip/shoulder/lumbar, pin joints at knee/elbow/neck, and
custom 2-DOF joints at ankle/wrist
(\texttt{OpenSim\_Models/src/opensim\_models/shared/body/body\_model.py}).
Every pack gives the pelvis an unconstrained 6-DOF root to the world ---
OpenSim's \texttt{FreeJoint} \texttt{ground\_pelvis}, MuJoCo's pelvis
\texttt{freejoint}, Drake's \texttt{floating} pelvis joint (configurable to
\texttt{fixed}), and Pinocchio's root built with
\texttt{pin.JointModelFreeFlyer()} --- except where a lift welds the pelvis to
a bench (below). The reference lifter is 80\,kg and 1.78\,m
(docs/development/lifting/PACK\_PARITY\_BASELINE.md, ``Setup''; Chapter 15,
``Conventions''); the bar is a 20\,kg competition bar plus 50\,kg per side,
120\,kg total.

Structural degrees of freedom measured at this commit
(docs/development/lifting/PACK\_PARITY\_BASELINE.md, ``Inventory''):

\begin{table}[h]
\centering\small
\begin{tabular}{lrrrr}
\toprule
Lift & MuJoCo $n_q/n_v$ & OpenSim $n_q/n_v$ & Drake $n_q/n_v$ & Pinocchio $n_q/n_v$\\
\midrule
Back Squat      & 56/52 & 34/34 & 35/34 & 35/34\\
Deadlift        & 56/52 & 34/34 & 35/34 & 35/34\\
Bench Press     & 56/52 & 28/28 & 28/28 & 35/34\\
Snatch          & 56/52 & 34/34 & 35/34 & 35/34\\
Clean and Jerk  & 56/52 & 34/34 & 35/34 & 35/34\\
\bottomrule
\end{tabular}
\caption{Generalized coordinate/velocity counts per pack and lift. Bench
press drops OpenSim and Drake to a welded (non-floating) root, which is the
``Root'' column of the same table.}
\end{table}

\subsection{Barbell and Plates}
Each pack models an Olympic barbell to IWF/IPF dimensions: a men's bar is
2.20\,m long, 1.31\,m between collars, 28\,mm shaft, 50\,mm sleeves, 20\,kg
unloaded; a women's bar is 2.01\,m, 25\,mm shaft, 15\,kg
(\texttt{OpenSim\_Models/src/opensim\_models/shared/barbell/barbell\_model.py},
\texttt{BarbellSpec}). The bar is three rigid bodies --- left sleeve, shaft,
right sleeve --- joined by welds, with plates added as sleeve mass via
\texttt{plate\_mass\_per\_side}. MuJoCo, Drake and Pinocchio each carry an
equivalent \texttt{BarbellSpec}/builder under their own \texttt{shared/barbell/}.
Bar and bench links are assumed to have zero mass-centre offset in every pack
at the recorded commits (Chapter 15, ``Assumptions and Limitations''); the lifter-only CoM
figure below depends on that assumption holding.

\subsection{Grip Interface}
The bar attaches to the hands with a rigid weld/fixed joint, not a compliant
contact, in every pack at this commit. OpenSim's \texttt{attach\_barbell\_to\_hands}
welds both hands to the shaft at a configurable \texttt{grip\_offset}
half-width; MuJoCo's \texttt{add\_weld\_constraint} does the same. Drake's
base builder welds the barbell to the left hand only (its own docstring calls
this ``SDF 1.8 kinematic-tree-safe'' for a single attachment point,
\texttt{Drake\_Models/src/drake\_models/exercises/base.py}); Pinocchio's
default \texttt{attach\_barbell} also welds via \texttt{add\_fixed\_joint} to
the left hand only. Measured attachment kinds per lift and pack are in the
``Hand and Bar Attachment'' table of
\texttt{docs/development/lifting/PACK\_PARITY\_BASELINE.md}; the right-hand
gap is in Section~\ref{sec:failed} below.

\subsection{Foot and Bench Contact}
OpenSim uses four Hunt-Crossley contact spheres per foot (heel/toe,
medial/lateral) against a ground half-space
(\texttt{OpenSim\_Models/src/opensim\_models/shared/body/foot\_contact.py}).
For the bench press it adds a near-zero-mass bench body welded to ground,
with the pelvis welded supine on it at the IPF-standard 0.43\,m bench height.
MuJoCo places a ground-plane geom in \texttt{worldbody} and uses its built-in
contact solver, with \texttt{<exclude>} pairs to suppress adjacent-segment
self-collision. Pinocchio defines contact frames and a Coulomb friction
coefficient for its own constraint-based contact --- a rectangular foot sole
(0.26\,m $\times$ 0.10\,m) with corner points --- rather than built-in
collision detection. Drake and Pinocchio expose no contact-force model at
all for the lift packs (Section~\ref{sec:gaps}, gap class \texttt{no\_grf}).
Contact-force parity should not be assumed from these geometry definitions
alone (\texttt{docs/help/movement\_optimizer.md}, ``Foot and Bench Contact'').

\subsection{Coordinate Frames and Units}
The audit's canonical frame is $+Z$ up, $+X$ forward, $+Y$ left, SI units.
MuJoCo, Drake and Pinocchio already use it; OpenSim is Y-up and its world
vectors are rotated into the canonical frame by a fixed rotation matrix
recorded in Chapter 15 (``Conventions''), which this document cites rather
than re-deriving. Positions are reported as body-frame origins in the world
frame, or pelvis-relative when comparing engines (Chapter 15, equation~4,
``placement-independent pose summary'').

\section{Phase Definitions and Trajectories}
Phase counts are not the same across packs, and in most lifts they are not
even the same within one engine versus the parity standard's own
\texttt{phase\_count} reference
(\texttt{docs/development/lifting/pack\_parity\_baseline.json},
\texttt{standard\_reference.phase\_count}: squat 5, deadlift 5, bench press 5,
snatch 6, clean and jerk 8).

\begin{table}[h]
\centering\small
\begin{tabular}{lrrrrr}
\toprule
Lift & Standard & MuJoCo & OpenSim & Drake & Pinocchio\\
\midrule
Back Squat      & 5 & 5 & 5 & 3 & 5\\
Deadlift        & 5 & 4 & 5 & 3 & 4\\
Bench Press     & 5 & 5 & 5 & 3 & 5\\
Snatch          & 6 & 6 & 6 & 5 & 6\\
Clean and Jerk  & 8 & 8 & 8 & 5 & 8\\
\bottomrule
\end{tabular}
\caption{Phase count per pack and lift against the standard's own reference
count (docs/development/lifting/PACK\_PARITY\_BASELINE.md, ``Inventory'' and
``Phases''). Drake is coarser everywhere (3 phases for squat/deadlift/bench,
5 for snatch/clean and jerk); MuJoCo and Pinocchio are short one deadlift
phase.}
\end{table}

Phase names differ per pack even where counts agree: MuJoCo's back-squat
phases are \texttt{descent\_start, half\_depth, bottom, drive, lockout};
OpenSim's are \texttt{start, descent, bottom, drive, lockout}; Drake's are
\texttt{unrack, bottom, lockout}; Pinocchio's are \texttt{start, descent,
bottom, drive, lockout}
(\texttt{docs/development/lifting/PACK\_PARITY\_BASELINE.md}, ``Phase names
per pack''). Pinocchio's phase targets are left-side only for every lift
(gap class \texttt{phase\_asymmetric}, Pinocchio\_Models\#449), and several
packs carry phase keys that the audit could not resolve to model coordinates
without guessing a side or coordinate (gap class \texttt{phase\_keys}:
OpenSim\_Models\#414, Drake\_Models\#390, Pinocchio\_Models\#436). MuJoCo's
phase values may be absolute joint positions or reference-relative angles;
the audit reads them as canonical angles, so MuJoCo phase out-of-range
findings are interpretation-dependent and are deliberately not used as gaps
(Chapter 15, ``Assumptions and Limitations'').

A shared phase standard across all four engines does not exist yet. LIFT-2
(\#11742) requires filing a \texttt{board-proposal} in
\texttt{Repository\_Management} first, for a single exercise specification
(anthropometric segment table, barbell spec, named phases with timing,
start/catch via-poses, root-joint policy, grip width/type per lift, foot
stance) owned in Tools and consumed by all four packs and by Movement
Optimizer (issue \#11742, ``Steps''). As of this revision that proposal is
tracked as board proposal \texttt{D-sorganization/Repository\_Management\#2060}
(Option A: extend the parity standard, schema
\texttt{biomech-parity-standard/v1}, to a v2.0 specification). The proposal
is open and awaits owner approval; no shared specification is implemented
before that decision.

\section{Inverse Dynamics, Spine Load and Muscle Redundancy}

\subsection{What Runs Today}
Three distinct things are easy to conflate here and are kept separate.
\emph{Inverse kinematics} against a physics-engine pack is not run: MuJoCo's
optimizer layer currently interpolates between named \texttt{Phase(name, t,
pose)} targets rather than solving IK, and switching to IK is an open,
unimplemented item (MuJoCo\_Models\#403). \emph{Feedback-assisted tracking} of
a lift controller is not run either: a controller API for the lift packs is
tracked as open (MuJoCo\_Models\#404) and no lift equivalent of the golf
same-input closed-loop controller (epic \#11605) exists. \emph{Independently
replayed open-loop dynamics} of a full lift is not run: the audit states
plainly that it ``does not simulate a lift'' (Chapter 15, ``Purpose''), and no
same-input bundle has been exported for any of the five lifts. All three are
therefore not run for the lift packs as of this revision --- not verified to
exist anywhere in the fleet for these models, and are reported here as absent
rather than guessed.

\subsection{Movement Optimizer}
Movement Optimizer (\texttt{Tools/src/movement\_optimizer}, vendored at
\texttt{src/shared/python/movement\_optimizer/}) is the one tool that produces
a full lift trajectory with inverse dynamics today, but on a different,
simpler body than the four packs above: a planar three-link shin/thigh/trunk
chain. Torques are recovered by Lagrangian inverse dynamics; the trajectory
itself is the result of multi-start parallel SLSQP under centre-of-mass
balance and joint-limit constraints, so \texttt{success} means the optimizer
converged and every hard constraint was satisfied, not that the trajectory is
globally optimal (\texttt{docs/help/movement\_optimizer.md}, ``Method'').
Exercises covered: \texttt{squat}, \texttt{full\_squat}, \texttt{deadlift},
\texttt{bench\_press}, \texttt{clean}, \texttt{jerk}, \texttt{snatch}. It is
explicitly not a physics-engine simulation and its results are not directly
comparable to the four engine packs' dashboards despite sharing a
\texttt{cross\_engine\_analysis} capability label
(\texttt{docs/help/movement\_optimizer.md}, ``Limitations'').

\subsection{Spine Load}
Movement Optimizer is the only source of a spine-load figure for the lifts.
\texttt{spine\_loads.py} computes the axial compression at the L5/S1 junction
as
\[
F_{\mathrm{comp}} = (m_{\mathrm{above}} + m_{\mathrm{bar}})\,g\,\cos(\theta_{\mathrm{torso}})
  + m_{\mathrm{torso}}\,\ddot\theta_{\mathrm{torso}}\,d_{\mathrm{torso}},
\]
with the L5/S1 joint placed at the base of the torso segment (the hip joint
in this three-link model), $m_{\mathrm{above}}$ the mass above L5/S1, and the
inertial term from the torso's own angular acceleration and its
centre-of-mass offset $d_{\mathrm{torso}}$
(\texttt{src/shared/python/movement\_optimizer/spine\_loads.py},
\texttt{spinal\_compression}). $m_{\mathrm{above}}$ depends on the exercise
family: for \texttt{squat}/\texttt{full\_squat} the torso segment already
includes the arms and \texttt{\_mass\_above\_l5} reads \texttt{body.m\_squat[2]};
for the deadlift family the torso is trunk plus head only and it reads
\texttt{body.m\_deadlift[2]}. The result is reported against
\texttt{NIOSH\_COMPRESSION\_LIMIT = 3400.0}\,N, the NIOSH recommended
occupational-lifting compression limit, which is an occupational-lifting
figure and not a clinical injury threshold for barbell training
(\texttt{spine\_loads.py}; \texttt{docs/help/movement\_optimizer.md},
``Limitations''). Only axial compression and anterior-posterior shear at one
level (L5/S1) are modelled; there is no lateral shear, no axial torsion, and
no other vertebral level. None of the four engine packs compute a spine load
at all.

\subsection{Muscle Redundancy}
Open, and depends on LIFT-6 (\#11746). None of the four engine packs
defines muscle actuators: the search recorded in
\texttt{docs/help/movement\_optimizer.md} (``What the Optimizer Computes'')
found no \texttt{Muscle}, \texttt{Thelen},
\texttt{Millard} or actuator classes in the OpenSim pack, consistent with the
epic's note that ``muscle content in the OpenSim pack is unverified'' (issue
\#11740). Movement Optimizer's own Hill-type muscle model
supplies torque-angle-velocity capacity and sticking-point detection
(\texttt{strength.py}) for its single three-link chain only; it does not
distribute load across named muscles and is not a muscle-redundancy solver
for a full-body musculoskeletal model. Full-body musculoskeletal lifts with
muscle forces, and muscle-redundancy solving across them, are future work
under LIFT-6.

\section{Measured Parity Results}
All figures below are read from the committed receipt,
\texttt{docs/development/lifting/pack\_parity\_baseline.json} (schema
\texttt{lift-pack-parity-baseline/v1}), and its rendered summary,
\texttt{docs/development/lifting/PACK\_PARITY\_BASELINE.md}, generated by
\texttt{scripts/lifting/run\_pack\_parity\_baseline.py} against pack commits
MuJoCo\_Models \texttt{735712cd63fe}, OpenSim\_Models \texttt{c854c4ed3256},
Drake\_Models \texttt{6ca7dcf18f78}, Pinocchio\_Models \texttt{6a4bb51a8fdb},
with parity standard version 1.2.0.

\begin{table}[h]
\centering\small
\resizebox{\textwidth}{!}{\begin{tabular}{lrrrrrr}
\toprule
Lift & Segments (m) & Hands (m) & Feet (m) & Bar centre (m) & Lifter CoM (m) & Mass spread (kg)\\
\midrule
Back Squat      & 0.0115 & 0.0115 & 0.0053 & 0.1000 & 0.0788 & 1.4e-05\\
Deadlift        & 0.0115 & 0.0115 & 0.0053 & 0.1915 & 0.0788 & 1.5e-05\\
Bench Press     & 0.0115 & 0.0115 & 0.0053 & 0.2406 & 0.0788 & 30\\
Snatch          & 0.0115 & 0.0115 & 0.0053 & 0.0200 & 0.0788 & 1.5e-05\\
Clean and Jerk  & 0.0115 & 0.0115 & 0.0053 & 0.1168 & 0.0788 & 1.5e-05\\
\bottomrule
\end{tabular}}
\caption{Worst pairwise same-input difference (m) over all engine pairs and
poses, pelvis-relative, against the 0.02\,m cross-engine position tolerance
(PACK\_PARITY\_BASELINE.md, ``Same-Input Cross-Engine Parity''). Segment,
hand and foot parity hold within tolerance for every lift (worst 0.0115\,m,
OpenSim vs.\ Drake at the flexion pose); bar-centre and lifter-CoM
differences exceed it everywhere. The bench-press mass spread (30\,kg) is the
bench-mass convention gap, not a body-segment discrepancy.}
\end{table}

Bar-centre differences come from the packs' own grip offsets, not from
forward kinematics (PACK\_PARITY\_BASELINE.md). Each engine's own FK was also
compared against the parity standard's reference FK, scaled to the lifter
height, worst over all lifts and poses: MuJoCo $2.22\times10^{-16}$\,m
(\texttt{foot\_r}); OpenSim $1.15\times10^{-2}$\,m (\texttt{forearm\_l});
Drake $3.33\times10^{-16}$\,m (\texttt{upper\_arm\_l}); Pinocchio
$2.22\times10^{-16}$\,m (\texttt{forearm\_l}) --- not evaluated for the bench
press (Chapter 15, equation~10). OpenSim's $1.15\times10^{-2}$\,m equals its
worst cross-engine segment difference above; its cause has not been
investigated.

Joint limits were checked against the standard for 28 coordinates per lift
and engine: OpenSim and Drake are off-standard on none of them in any lift;
MuJoCo is off on 4--8 coordinates depending on the lift (hip/knee/shoulder/wrist,
up to 1.396\,rad deviation on hip flexion); Pinocchio is off on 10
coordinates in every lift (ankle, lumbar and neck, up to 0.698\,rad on neck
flexion) (PACK\_PARITY\_BASELINE.md, ``Joint Limits Versus the Standard'').

\section{Discrepancies and Open Gaps}\label{sec:gaps}
Every discrepancy below is computed from the measured rows by
\texttt{derive\_gaps} (\texttt{src/shared/python/lifting/pack\_audit/gaps.py});
none is entered by hand. Each maps to per-pack tracking issues and to the
LIFT child of epic \#11740 expected to close it
(PACK\_PARITY\_BASELINE.md, ``Discrepancies and Tracking Issues'').

\begin{table}[h]
\centering\footnotesize
\begin{tabular}{p{4.6cm}lp{3.6cm}l}
\toprule
Discrepancy & Engines & Example issues & Epic child\\
\midrule
Right hand not attached to the bar & Drake, Pinocchio & Drake\_Models\#365; Pinocchio\_Models\#443 & LIFT-3\\
Grip half-width vs.\ achievable spacing & Drake, MuJoCo, OpenSim, Pinocchio & MuJoCo\_Models\#408; OpenSim\_Models\#394 & LIFT-3\\
Floor-pull start pose not at the floor & Drake, MuJoCo, OpenSim, Pinocchio & Drake\_Models\#381, \#364; MuJoCo\_Models\#407 & LIFT-4\\
Joint limits differ from the standard & MuJoCo, Pinocchio & MuJoCo\_Models\#407, \#409; Pinocchio\_Models\#433 & LIFT-5\\
Phase count differs from the standard & Drake, MuJoCo, Pinocchio & Drake\_Models\#372; MuJoCo\_Models\#401 & LIFT-6\\
Phase targets are left-side only & Pinocchio & Pinocchio\_Models\#449 & LIFT-6\\
Phase keys don't resolve to model coords & Drake, OpenSim, Pinocchio & Drake\_Models\#390; OpenSim\_Models\#414 & LIFT-6\\
Bench mass convention differs & Drake, MuJoCo, OpenSim, Pinocchio & Drake\_Models\#391; MuJoCo\_Models\#428 & LIFT-2\\
Start pose carries unphysical contact force & MuJoCo, OpenSim & MuJoCo\_Models\#427; OpenSim\_Models\#384 & LIFT-4\\
No ground-reaction output available & Drake, Pinocchio & Drake\_Models\#371; Pinocchio\_Models\#432 & LIFT-9\\
Segment inertia differs from the other packs & Drake & Drake\_Models\#391 & LIFT-2\\
\bottomrule
\end{tabular}
\caption{All eleven discrepancy classes derived by \texttt{derive\_gaps} at
this baseline. ``Example issues'' abbreviates each row; the full issue list
per class is in PACK\_PARITY\_BASELINE.md.}
\end{table}

\section{Failed Experiments and Negative Results}\label{sec:failed}
These are the negative results actually documented for the lift packs. No
additional failed-experiment narratives (in the style of, for example, the
same-input-parity reference's ``Failed First Comparison''/``Failed First
Version'' sections) exist in the lift-pack sources read for this document;
none are synthesized here.

\begin{itemize}[nosep]
\item \textbf{No floor-pull start pose reaches the plate radius.} The audit
  threshold is $0.225\pm0.05$\,m (450\,mm plate radius; Chapter 15,
  ``Tolerances''). Measured bar-centre-above-sole at the deadlift start pose:
  MuJoCo 1.017\,m, OpenSim 0.7097\,m, Drake 0.6969\,m, Pinocchio 0.6586\,m ---
  every pack, every floor-pull lift (deadlift, snatch, clean and jerk), misses
  the target (PACK\_PARITY\_BASELINE.md, ``Start Poses''; gap class
  \texttt{start\_pose}).
\item \textbf{The right hand is not attached to the bar in Drake or
  Pinocchio.} Both weld/fix only the left hand; Pinocchio's right-hand
  attachment is only a dummy frame, \texttt{barbell\_grip\_r}
  (PACK\_PARITY\_BASELINE.md, ``Hand and Bar Attachment''; gap class
  \texttt{grip\_attachment}).
\item \textbf{Grip half-width does not match achievable hand spacing in any
  pack.} For the deadlift, intended half-width vs.\ achievable half-spacing:
  MuJoCo 0.22\,m vs.\ 0.1709\,m (mismatch 0.0491\,m); OpenSim 0.4\,m vs.\
  0.1594\,m (mismatch 0.2406\,m); Drake 0.22\,m vs.\ 0.1709\,m (mismatch
  0.0491\,m); Pinocchio 0.393\,m vs.\ 0.1709\,m (mismatch 0.2221\,m), all
  above the 0.01\,m audit threshold (PACK\_PARITY\_BASELINE.md, ``Hand and Bar
  Attachment''; gap class \texttt{grip\_width}).
\item \textbf{The bench-press start pose carries an unphysical contact
  force in OpenSim.} 21.7\,kN of vertical ground contact force at the bench
  default pose (PACK\_PARITY\_BASELINE.md, ``Epic Defect List Review'';
  confirmed finding for OpenSim\_Models\#384). MuJoCo's own start poses carry
  large non-ground contact force for the floor-pull lifts (up to
  $4.486\times10^{5}$\,N on the snatch), a separate instance of the same gap
  class \texttt{contacts}.
\end{itemize}

\section{Limitations}
\begin{itemize}[nosep]
\item \textbf{Unavailable is never zero.} Drake and Pinocchio expose no
  contact model, so start-pose ground-reaction force is reported
  unavailable with a reason string, never as 0\,N (Chapter 15, ``Unavailable
  Values''; gap class \texttt{no\_grf}).
\item Native hand-to-bar closure is reported only where an engine defines it
  (MuJoCo welds, OpenSim weld constraint after assembly); it is unavailable
  for Drake and Pinocchio.
\item The lifter-only centre of mass is unavailable when no lifter mass
  remains, and bar height above the sole is unavailable when no foot-contact
  geometry is found. Missing bodies, non-finite positions and coincident bar
  sleeves raise \texttt{ValueError}; no geometric value is filled in.
\item The bar/bench zero-mass-centre-offset assumption (Section~2.2) holds at
  the recorded commits; the lifter-only CoM figure is wrong if a pack changes
  that without updating the audit.
\item OpenSim's start-pose coordinate values are read after model assembly,
  which moves coordinates off their declared defaults.
\item OpenSim, Drake and Pinocchio stills are skeleton plots of each engine's
  own FK, not native renders.
\item No Simscape or MATLAB lift models exist in any repository, and nothing
  was substituted for them (epic \#11740).
\item This document is a software-consistency reference. Scientific
  qualification of any lift model against measured data stays with the
  design-manual governance pathway (\texttt{scripts.check\_design\_manual\_governance});
  a successful build of this document is not scientific, semantic, visual,
  accessibility, or publication approval of anything.
\end{itemize}

\section{Reproduction}
From the repository root, with the four pack checkouts found under
\texttt{LIFT\_PACK\_ROOT} (Chapter 15, ``Tests and Reproduction''):
\begin{verbatim}
export PYTHONPATH=.:src MUJOCO_GL=egl MPLBACKEND=Agg QT_QPA_PLATFORM=offscreen
python3 scripts/lifting/run_pack_parity_baseline.py
python3 scripts/lifting/render_pack_stills.py   # optional stills
\end{verbatim}
This rewrites \texttt{docs/development/lifting/pack\_parity\_baseline.json}
and its rendered Markdown. The geometry, frame and gap-derivation tests run
without any engine installed:
\begin{verbatim}
python3 -m pytest tests/unit/lifting/pack_audit/test_pack_audit_geometry.py \
  tests/unit/lifting/pack_audit/test_pack_audit_frames_names.py \
  tests/unit/lifting/pack_audit/test_pack_audit_report.py -q
\end{verbatim}
\texttt{test\_pack\_audit\_live.py} runs the adapters against checked-out
packs and skips when the packs or engines are absent. This document itself
compiles with:
\begin{verbatim}
pdflatex -interaction=nonstopmode -halt-on-error lift_models.tex
\end{verbatim}
run twice to resolve references; the generated PDF is not committed.

\end{document}
